\documentclass{beamer}

\newcommand{\week}{Banks 3 of 6}

\title{Bank Earnings}
\subtitle{Mishkin Chapter 9}
\author{Econ 357}
\date{\week}

% Reference the shared preamble
\input{preamble}

% Set the footer -- change 
\setbeamertemplate{footline}{
  \leavevmode%
  \vspace{2ex}
  \hbox{%
    % Left box: Econ 457
    \begin{beamercolorbox}[wd=.4\paperwidth,ht=2.5ex,dp=1ex,left]{author in head/foot}%
      \hspace{1em}Econ 357
    \end{beamercolorbox}%
    % Middle box: Week
    \begin{beamercolorbox}[wd=.2\paperwidth,ht=2.5ex,dp=1ex,center]{date in head/foot}%
      \centering\week
    \end{beamercolorbox}%
    % Right box: Slide numbers
    \begin{beamercolorbox}[wd=.4\paperwidth,ht=2.5ex,dp=1ex,center]{date in head/foot}%
      \hfill\insertframenumber{} 
    \end{beamercolorbox}%
  }%
  \vskip0pt%
}

\begin{document}

\frame{\titlepage}

\begin{frame}
    \frametitle{Outline for Banks}
        \begin{enumerate}
            \item Role of Banks
            \item Bank Balance Sheets
            \item \fbox{Bank Earnings}
            \item Bank Management (and Investment Banks)
            \item Bank Failures
            \item Bank Regulation
        \end{enumerate}
\end{frame}


\begin{frame}
    \frametitle{Outline for Today's Lecture}
    
    \begin{enumerate}
        \item Review: Bank Balance Sheets
        \item Required Ratios
        \item Balance Sheet Changes
        \begin{itemize}
        \item Practice
        \end{itemize}
        \item How Banks Make Money
        \begin{itemize}
            \item Net Interest Margin
            \item Impact of Leverage
            \item Practice
        \end{itemize}
    \end{enumerate}

\end{frame}

\begin{frame}[t]
    \frametitle{1. Review: Bank Balance Sheets}
    \framesubtitle{Assets}

    \centering
    \includegraphics[width=0.8\textwidth]{Econ357/figures/03_banks_02_assets.png}

    \blfootnote{Source: FDIC Quarterly Banking Profile, 2024-q2}
    
\end{frame}

\begin{frame}[t]
    \frametitle{1. Review: Bank Balance Sheets}
    \framesubtitle{Liabilities}

    \centering
    \includegraphics[width=0.8\textwidth]{Econ357/figures/03_banks_02_liabilities.png}

    \blfootnote{Source: FDIC Quarterly Banking Profile, 2024-q2}
    
\end{frame}

\begin{frame}[t]
    \frametitle{1. Review: Bank Balance Sheets}
    \framesubtitle{Capital}

    A hypothetical bank's capital stack may look like this:

    \centering
    \includegraphics[width=\textwidth]{Econ357/figures/03_banks_02_bank_capital_stack.png}

\end{frame}

\begin{frame}[t]
    \frametitle{2. Required Ratios}

    There are two important limits that banks must adhere to:
    \vspace{1em}

    \begin{enumerate}
        \item \textbf{Reserve Requirements}: Banks are required to maintain a minimum amount of reserves.   Usually reserves must be maintained at $\approx$10\% deposits.
        $$\text{Required Reserves} = \frac{\text{Reserves}}{\text{Deposits}}$$
        \item \textbf{Capital Ratios}: Banks are required to maintain a minimum amount of capital.   (You could alternatively think about this as a restriction on bank leverage)  The minimum capital ratio for most banks is around 10\% and could be as high as 15\% for larger banks.
        $$\text{Capital Ratio} = \frac{\text{Capital}}{\text{Assets}}$$
    \end{enumerate}
    
\end{frame}

\begin{frame}[t]
    \frametitle{3. Bank Balance Sheet Changes}

    What happens when you deposit money at a bank?
    
    \begin{table}
    \centering
    \caption{\textbf{Your Bank}}
    \begin{tblr}{
      colspec = { Q[c,wd=2cm] Q[c,wd=1cm] |Q[c,wd=2cm] Q[c,wd=1cm] }, vline{3} = {1-Z}{}
    }
    \toprule
    Assets && Liabilities \\
    \midrule
    Cash & +\$100 & Deposits & +\$100 \\
    \end{tblr}
    \end{table}

\end{frame}

\begin{frame}[t]
    \frametitle{3. Bank Balance Sheet Changes}

    What happens when you write a check to somebody else?

    \begin{table}
    \centering
    \caption{\textbf{Your Bank}}
    \begin{tblr}{
      colspec = { Q[c,wd=2cm] Q[c,wd=1cm] |Q[c,wd=2cm] Q[c,wd=1cm] }, vline{3} = {1-Z}{}
    }
    \toprule
    Assets && Liabilities \\
    \midrule
    Cash & -\$100 & Deposits & -\$100 \\
    \end{tblr}
    \end{table}

        \begin{table}
    \centering
    \caption{\textbf{Recipients' Bank}}
    \begin{tblr}{
      colspec = { Q[l,wd=2cm] Q[c,wd=1cm] |Q[l,wd=2cm] Q[c,wd=1cm] }, vline{3} = {1-Z}{}
    }
    \toprule
    Assets && Liabilities \\
    \midrule
    Cash & +\$100 & Deposits & +\$100 \\
    \end{tblr}
    \end{table}

\end{frame}

\begin{frame}[t]
    \frametitle{3. Bank Balance Sheet Changes}

    What happens you (“A”) deposit money at a bank and the bank then makes a loan to somebody else (“B”) that uses the same bank?   

    {\small
    \begin{table}
    \centering
    \caption{\textbf{Step 1}}
    \begin{tblr}{
      colspec = { Q[l,wd=3cm] Q[c,wd=1cm] |Q[l,wd=3cm] Q[c,wd=1cm] }, vline{3} = {1-Z}{}
    }
    \toprule
    Assets && Liabilities \\
    \midrule
    Cash & +\$100 & Deposit from A & +\$100 \\
    \end{tblr}
    \end{table}
    \vspace{-1em}

        \begin{table}
    \centering
    \caption{\textbf{Step 2}}
    \begin{tblr}{
      colspec = { Q[l,wd=3cm] Q[c,wd=1cm] |Q[l,wd=3cm] Q[c,wd=1cm] }, vline{3} = {1-Z}{}
    }
    \toprule
    Assets && Liabilities \\
    \midrule
    Loan to B & +\$100 & Deposit from A & +\$100 \\
    Cash & +\$100 & Deposit from B & +\$100 \\
    \end{tblr}
    \end{table}
    }

\end{frame}

\begin{frame}[t]
    \frametitle{3. Bank Balance Sheet Changes}

    In the previous example the bank literally created money.   It received a \$100 deposit, but ended up with \$200 in deposits.\\
    \vspace{1em}
    Required reserves limit this process and slow down money creation.\\
    \vspace{1em}
    To see how, the two next slide shows the same process but now with the requirement that the bank hold reserves equal to at least 10\% of deposits and then with capital ratio requirement.   In both case, the bank is limited in how much lending it can do, which slows down the money creation process.  

\end{frame}

\begin{frame}[t]
    \frametitle{3. Bank Balance Sheet Changes}

    Assume the bank must maintain required reserves of at least 10\% of deposits.

    {\small
    \begin{table}
    \centering

    \begin{tblr}{
      colspec = { Q[l,wd=3cm] Q[c,wd=1cm] |Q[l,wd=3cm] Q[c,wd=1cm] }, vline{3} = {1-Z}{}
    }
    \toprule
    Assets && Liabilities \\
    \midrule
    Cash & +\$90 & Deposit from A & +\$100 \\
    Reserves & +\$10 & && \\
    \end{tblr}
    \end{table}

    \begin{table}
    \centering
    \begin{tblr}{
      colspec = { Q[l,wd=3cm] Q[c,wd=1cm] |Q[l,wd=3cm] Q[c,wd=1cm] }, vline{3} = {1-Z}{}
    }
    \toprule
    Assets && Liabilities \\
    \midrule
    Loan to B & +\$90 & Deposit from A & +\$100 \\
    Cash & +\$81 & Deposit from B & +\$90 \\
    Reserves & +\$19
    \end{tblr}
    \end{table}
    }
    Because of the required reserves, the bank only created \$90 (instead of \$100 in the first example).

\end{frame}

\begin{frame}[t]
    \frametitle{3. Bank Balance Sheet Changes}

    Assume the bank starts with \$15 in capital and must maintain a capital ratio of at least 10\%.


    {\small
    \begin{table}
    \centering
    \begin{tblr}{
      colspec = { Q[l,wd=3cm] Q[c,wd=1cm] |Q[l,wd=3cm] Q[c,wd=1cm] }, vline{3} = {1-Z}{}
    }
    \toprule
    Assets && Liabilities \\
    \midrule
    Cash & +\$90 & Deposit from A & +\$100 \\
    Reserves & +\$10 & && \\
    && Capital & +\$15\\
    \end{tblr}
    \end{table}

    \begin{table}
    \centering
    \begin{tblr}{
      colspec = { Q[l,wd=3cm] Q[c,wd=1cm] |Q[l,wd=3cm] Q[c,wd=1cm] }, vline{3} = {1-Z}{}
    }
    \toprule
    Assets && Liabilities \\
    \midrule
    Loan to B & +\$50 & Deposit from A & +\$100 \\
    Cash & +\$85 & Deposit from B & +\$50 \\
    Reserves & +\$15 \\
        && Capital & +\$15\\

    \end{tblr}
    \end{table}
    }

\end{frame}


\begin{frame}[t]
    \frametitle{3. Bank Balance Sheet Changes}

    A couple more examples:
    \vspace{1em}
    
    Deposit outflows:
    \begin{itemize}
        \item Use most liquid assets
        \item Need to maintain reserve requirement
    \end{itemize}
    
    Loan Losses
    \begin{itemize}
        \item Regulators may force raising more capital
        \item This involves getting new owners to contribute cash in exchange for equity. What are the pros/cons of that?
        \item If losses are large enough, may have negative capital (will discuss this more in a later class...)
    \end{itemize}

\end{frame}

\begin{practiceframe}[t]
    \frametitle{3. Bank Balance Sheet Changes}
    \vspace{-1em}
    Consider a bank with the following: \$30 million in capital, \$20 million in debt, \$200 million in deposits, \$50 million in securities, \$150 million in loans, and the remainder of its assets in cash/reserves. The banks regulators require it to maintain a capital to asset ratio of 10\% and a reserve to deposit ratio of 10\%.
    \vspace{1em}
    \begin{enumerate}
        \item Write down the bank’s balance sheet using a “T-chart”. Verify that it is meeting both of its regulatory ratios.
        \item What happens if there is a deposit outflow of \$40 million? Write down the bank balance sheet. Is the bank currently meeting its regulatory ratios?
        \item What happens if instead of a deposit outflow,the bank experienced losses on its loans. Specifically, what happens if 10\% of the banks’ loans default? Write down the bank balance sheet.
        \item Is the bank currently meeting its regulatory ratios?
    \end{enumerate}
    
\end{practiceframe}

\begin{frame}[t]
    \frametitle{4. How Banks Make Money}
    \framesubtitle{Net Interest Income}

    \centering
    \textit{Net Interest Income = Return on Assets - Cost of Funding}
    \vspace{1em}

    \begin{itemize}
        \item Bank Assets = Loans, securities, cash, reserves
        \item Bank Liabilities (funding) = Deposits, borrowing
    \end{itemize}
    \vspace{1em}

    \centering
    \includegraphics[width=0.7\textwidth]{Econ357/figures/03_banks_3_net_interest.png}

    \blfootnote{FDIC Quarterly Banking Profile, Q2-2024}

\end{frame}

\begin{frame}[t]
    \frametitle{4. How Banks Make Money}
    \framesubtitle{Net Interest Income}

    \centering
    \includegraphics[width=0.8\textwidth]{Econ357/figures/03_banks_3_nii_chart.png}

    \blfootnote{FDIC Quarterly Banking Profile, Q2-2024}

\end{frame}

\begin{frame}[t]
    \frametitle{4. How Banks Make Money}
    \framesubtitle{Impact of Leverage}

\begin{itemize}
    \item "ROA" stands for the "Return on Assets."   ROA is typically expressed as a percent, similar to the total return on a financial security.
    \item ROA = Net Profit / Total Assets
    \item "ROE" stands for "Return on Equity."   ROE is also typically expressed as a precent.
    \item ROE = Net Profit / Equity Capital
\end{itemize}
\vspace{1em}

\centering
\fbox{ROE = ROA X (Assets / Equity Capital)}

\end{frame}

\begin{frame}[t]
    \frametitle{4. How Banks Make Money}
    \framesubtitle{Impact of Leverage}

    \centering
    \fbox{ROE = ROA X (Assets / Equity Capital)}
    \vspace{1em}

    \begin{itemize}
        \item Assets/Equity Capital = ”Equity Multiple” (also the leverage ratio, which we covered two lectures ago)
        \item This is the inverse of the bank’s “capital ratio”
        \item Note that a higher “Equity Multiple” (or, equivalently a lower “capital ratio”) increases the ROE
        \item Bank management therefore has an incentive to have lower capital and higher ROE.   
        \item Why may this not be desirable for everybody else?
    \end{itemize}

\end{frame}

\begin{frame}[t]
    \frametitle{4. How Banks Make Money}
    \framesubtitle{Impact of Leverage}

     \centering
    \textit{Net Income = Net Interest Income - Taxes - Cost of Bad Loans + Securities Gains - Cost of Doing Business}
    \vspace{1em}

    \centering
    \includegraphics[width=0.8\textwidth]{Econ357/figures/03_Banks_3_income.png}

    \blfootnote{FDIC Quarterly Banking Profile, Q-2024}
    
\end{frame}

\begin{frame}[t]
    \frametitle{4. How Banks Make Money}
    \framesubtitle{Impact of Leverage}

    \centering
    \includegraphics[width=0.8\textwidth]{Econ357/figures/03_Banks_3_income2.png}

\end{frame}

\begin{practiceframe}[t]
    \frametitle{4. How Banks Make Money}

    Consider the same bank as before: \$30 million in capital, \$20 million in debt, \$200 million in deposits, \$50 million in securities, \$150 million in loans, and the remainder of its assets in cash/reserves. 
    \vspace{1em}

    Assume the following returns on assets: cash = 0\%, securities = 5\%, loans = 10\%.  Also assume the follow cost of funding: deposits = 1\%, debt = 4\%.  What is the banks’ return on assets?  What is the banks’ return on equity? 
    \vspace{1em}
    
    \begin{itemize}
        \item What is the bank’s return on assets?
        \item What is the bank’s return on equity?
    \end{itemize}
    
\end{practiceframe}

\end{document}